\documentclass[script,pro]{debate}
\motion{语言的边界是 / 不是人类的边界}
\stance{正方}{是}
\meta{赛制}{招新赛 3v3}
\meta{口径来源}{备赛文档 第十节 正方口径表}
\begin{document}
\debatetitle
\begin{thesisbox}[全场一句话立论]人比别的动物多走出去的每一步，都是语言铺的路。语言到不了的地方，一个人也许摸得到，人类走不到。\end{thesisbox}

\begin{speech}{1. 正一开篇立论稿}{正文 824 字 / 180 秒}
谢谢主席，各位好。

维特根斯坦有一句名言：我的语言的界限，意味着我的世界的界限。今天的辩题把“我的”换成了“人类的”。一字之差，问题就从一个人能摸到哪里，变成了人类能走到哪里。

先说清三个词。语言，是有词、有语法、用来表达和交流意思的符号系统，口语、文字、手语都是；音乐、绘画叫语言，只是比喻。人类的边界，是人能想到、认识到、做到，并且能传给别人的范围的极限。“是”，指两条边界基本重合。

所以判断本题的标准是：语言到不了的地方，人类还能不能稳定地到达。正方要证明，人超出本能的那一截以语言为条件，语言之外的，要么到不了，要么留不住；反方要证明，在语言之外，人能稳定地思考、创造，而且留得住。只说人有说不清的时候，不够。

我方第一点：没有词，走不到。

人天生有一些底子：三以内的数，大概多少，眼前的东西。很多动物也有。人比动物多出来的那一截，精确的数，明白别人心里想错了什么，是跟着语言一起到的。

我举个例子。尼加拉瓜的聋校里，几十年前出现了一门全新的手语。最早学它的那批人，学到的版本几乎没有“想”“知道”这类词。心理学家给他们做过一道题：一个人出门时东西放在这儿，有人趁他不在挪走了，他回来会去哪儿找？第一批八个人，平均二十六七岁，四道题平均只答对半道。后一批学到了这些词的年轻人，将近四道全对。作者的原话是：二十五年的社会经验，也替代不了语言。

维果茨基说，思想是在词里完成的。语言走到哪，人才走到哪。

我方第二点：说不出就接不上。

人类和别的动物最大的不同，是能一代接着一代往前走。可接力要传得准。二〇一五年有个实验，一百八十四个人排成传递链，学打最古老的奥杜威石器。只看师傅打、不能交流的，每一下敲出好石片的概率，没有提高。只许比划的，翻了一倍。能说话的，翻了四倍。而历史上的奥杜威石器，七十多万年几乎没变样，作者推测，这和传得不准有关。

庄子讲过一个老木匠，叫轮扁。他说做车轮的手艺“口不能言”，所以“臣不能以喻臣之子”，儿子学不走。下一句更重：古之人与其不可传也死矣。说不出的东西，跟着一个人活，也跟着一个人走。

说白了，一个人摸到了，那是他的边界。说出来，传下去，人类才走到。

所以我方全场会请对方回答一个问题：\stress{说不出的东西，怎么让下一个人接着往前走？}

我方承认，人有说不清的时刻。但人类走到今天，靠的是有人把它说清楚。谢谢。
\end{speech}
\begin{contingency}
\ifthen{对方把语言定义成“嘴上说的话”}{手语是不是语言？聋人有没有语言？把手语赶出去，是在用定义取消聋人，请回到中立定义。}
\ifthen{对方举失语的人照样思考}{先承认能，再说脚手架：他们都是先有语言、后来才失去的；一辈子没有语言的人，才是检验。}
\ifthen{对方说 5.6 讲的是“我的”不是“人类的”}{同意，所以我方论证的是人类，用的是实验，不是名言。}
\ifthen{对方用“言不尽意”}{说不尽的那一半，留在了一个人那里，这正是我方说的边界。}
\end{contingency}

\begin{stage}{2. 正一被质询防守预案}{反二质询，90 秒双边计时}
\textbf{原则}：先答是或否，再用一句话守住前提。不否认常识：失语的人能思考，婴儿会算小数目，都先认。对方打断时停下来，等问完再答；每答不超过十秒。
\begin{qa}
\qarow{失语的人还能不能思考？}{能。他们的思考是有语言的几十年里搭起来的；楼盖好了拆掉脚手架，楼还在。}
\qarow{他的世界有没有跟着他的语言变小？}{他一个人的世界没变小；人类能接到他那里的东西，停在了他说出口的最后一句。}
\qarow{你方是不是说没有语言就不能思考？}{不是。底子人人有，很多动物也有；我方说的是人超出动物的那一截。}
\qarow{尼加拉瓜手语是不是那群聋孩子自己造的？}{是。在造出“想”“知道”之前，他们走不到那一步；边界一起挪，挪的办法是造词。}
\qarow{只比划不出声，是不是也翻了一倍？}{是。比划帮人守住，一倍；说话让人往前走，四倍。}
\qarow{婴儿五个月会算一加一，算不算语言之外的思考？}{算，那是底子，猴子身上也测到过。婴儿要等学会数词，才到得了十七。}
\qarow{爱因斯坦说他思考时不用词？}{他先摸到；他自己说要到第二阶段去找词，为的是能告诉别人。找到了，人类才走到。}
\qarow{维特根斯坦自己也说有不可说的？}{说了，他让我们对它沉默。沉默的东西，一个人可以活出来，人类没法接着往前推。}
\end{qa}
\end{stage}

\begin{stage}{3. 正二质询反一问题链}{90 秒，双边计时}
\textbf{总目标}：拿到两个承认。一，对方举的失语者，思考是在有语言的时候搭起来的；二，说不出的那一点，每一代都要重新摸。问题要短，一句话答得完；对方绕开时先复述一次问题，第二次再礼貌打断。
\begin{chain}{链一：脚手架}{失语者是先有语言的}
\q{您方举的失语者，发病前会不会说话？}{答“会”：下一问；答“不一定”：请举一个从没学过语言、却会解括号算式的人}
\q{他们的算术，是不是在会说话的时候学会的？}{答“是”：下一问；答“否”：请问是在哪里学会的}
\q{一辈子没有语言的人，您方有能精确数过三的例子吗？}{答“没有”：收束；答“有”：请报出处，我方会在申论里核对}
\cut{好，我换个问法：他们学算术的时候，会不会说话？}
\close{谢谢。对方承认，失语者的思考是在有语言的时候搭起来的。}
\end{chain}
\begin{chain}{链二：接力}{说不出的每代重摸}
\q{您方说的手感，是不是说不出来？}{答“是”：下一问}
\q{那说不出的那一点，是不是每一代都要重新摸？}{答“是”：下一问；答“否”：请问怎么传}
\q{师傅走了，徒弟没摸到，人类还有没有这一点？}{答“没有”：收束；答“有”：请问在哪里}
\close{谢谢。说不出的那一点，每一代从零开始。请评委记住：守得住，接不上。}
\end{chain}
\end{stage}

\begin{stage}{4. 正三对辩要点}{双方各 90 秒}
\textbf{目标}：开第一个战场，思考。让评委看到：对方所有“语言之外的思考”，要么是先有语言的人，要么是动物也有的底子。每次发言先一句回应，再抛一个问题。
\begin{points}{进攻点（4–5 个，每个 15–20 秒）}
\item 失语的人都是先学会语言的。请给我一个一辈子没有语言、却能精确数到十的人。
\item 尼加拉瓜第一批八个人，二十五年的社会经验，四道题只对半道。他们缺的是经验，还是那几个词？
\item 婴儿会的一加一，猴子身上也测到过。那是人类的边界，还是动物的边界？
\item 您方的例子都是一个人摸到了。这个人走了以后，人类还剩下什么？
\item 您方说语言只是工具。那请问，没有这个工具，人类能不能把十七这个数记过一夜？
\end{points}
\begin{points}{备用点（6–8 个）}
\item 对方引 Fedorenko：同一篇写着，语言是传递文化知识的强大工具。
\item 对方引爱因斯坦：他说要到第二阶段去找词，为的是告诉别人。
\item 对方说造手语的是人：人造词，词带人走，这叫一起挪。
\item 对方说默会知识：一个人知道的比说的多，人类知道的不比说的多。
\item 对方说比划也翻一倍：一倍是守住，四倍是往前走。
\item 对方说皮拉罕人摆在眼前能配对：东西一拿走就不准了，那就是边界。
\item 对方说言不尽意：说不尽的那一半，就留在了一个人那里。
\end{points}
\begin{points}{对方最可能问的三个问题 → 回应}
\item 失语的人还能不能思考？→ 能，是语言先搭起来的；而且他想的都是语言说得清的事。
\item 你说不出你妈长什么样，你认不认得她？→ 认得。认脸很多动物也会，那不是人类的边界。
\item 尼加拉瓜那门手语是谁造的？→ 人造的；造出来之前，他们走不到。
\end{points}
\end{stage}

\begin{kit}{5. 正二申论}{套件 · 组装后 385–402 字 / 90 秒}
\assembly{开头 → 回应反二（任选 1） → 拆反方立论（任选 1） → 推进论点一 → 收尾}
\begin{fixed}{开头}
谢谢主席。我先说这一篇做什么：拆开对方的立论，再往前推我方第一点。
\end{fixed}
\begin{pick}{回应反二}{1}
\begin{module}{如果反二申论主打失语者还能思考}
反二刚才讲，失语的人照样能思考。我方从头到尾都承认，他们能想。可他们的思考，是在有语言的几十年里一点点搭起来的。楼盖好了，拆掉脚手架，楼还在；这证明不了楼不靠脚手架盖。
\end{module}
\begin{module}{如果反二申论主打爱因斯坦先想后找词}
反二刚才用了爱因斯坦。请评委听听那封信的后半句：词语要到第二阶段才去找，为的是能告诉别人。他一个人先摸到了；人类是等他把词找到、把论文写出来，才走到的。
\end{module}
\begin{module}{如果反二申论主打石器没有语言也传了七十万年}
反二刚才说，石器没有像样的语言，也传了七十多万年。我方同意，传了。可这七十多万年，它几乎没变样。守住和往前走，是两回事；人类的边界，是一代一代往前走出来的。
\end{module}
\begin{fallback}{反二申论不在以上几条时}
对方刚才那一篇，请评委用判准量一量：它讲的是一个人摸到了什么，还是人类稳定地走到了哪里。这两件事之间，差着的正是今天这一整道题，也是双方真正的分歧。
\end{fallback}
\end{pick}
\begin{pick}{拆反方立论}{1}
\begin{module}{如果对方立论建立在失语者还能思考上}
对方立论的主干，是失语的人还能思考。可请看他们想的是什么：算术、逻辑题、下棋。这些事，语言说得清清楚楚，教科书里一条条写着。一个人把说话的工具弄丢了，不等于他走到了语言到不了的地方。对方要证明的，是边界外面还有人；拿出来的，却是边界里面的人，只是嘴坏了。说白了，这是工具出了故障，不是边界的反例。
\end{module}
\begin{module}{如果对方立论建立在默会知识上}
对方立论的主干，是我们知道的比说出来的多。我方承认，一个人知道的，比他说得出的多。可人类知道的，不比说出来的多。轮扁的故事，大家常常只记前一半，下一句是：古之人与其不可传也死矣。说不出的那一点，徒弟只能从零去摸，摸不到就断了，下一代再从零开始。一代一代从零开始，这正是我方说的边界。
\end{module}
\begin{module}{如果对方立论建立在语言只是交流工具上}
对方立论的主干，是语言只是交流的工具，不是思考的材料。好，就按对方说的，它是工具。可提出这个说法的那篇《自然》综述，同一篇摘要里还写着：语言是传递文化知识的强大工具。人类不是一个人，是一代接一代的人。传递的工具到不了哪里，下一代就接不到哪里。工具的边界，就是接力的边界，也就是人类的边界。
\end{module}
\begin{fallback}{对方立论的主干不在以上几条时}
我方想请评委回到判准。一个人的例子举得再多，也还停在一个人身上。对方要证明的，不是人有说不清的时候，而是在语言之外，人能稳定地到达，还留得住。一个说不清的瞬间，证明不了边界；一个人摸到的东西，也证明不了人类走到。这两步，只要对方还没有走完，“是”这个字就还站在我方这一边。
\end{fallback}
\end{pick}
\begin{fixed}{推进论点一}
我再往前推一步。尼加拉瓜还有一群更孤单的人，叫家庭手势者。他们一辈子没接触过任何一门语言，只在家里比划。可他们有工作，会挣钱，认得钱。二〇一一年有个研究，找了四个这样的成年人。一到三个东西，他们全对。一超过三，数得接近，精确答对的不到一半。你想想，他们不缺社会，不缺经验，缺的是那几个数词。有了词，才数得到十七。
\end{fixed}
\begin{fixed}{收尾}
所以请评委记住我方的问题，也请对方回答：说不出的东西，怎么让下一个人接着往前走？
\end{fixed}
\end{kit}

\begin{stage}{6. 正二被质询防守预案}{反三质询，90 秒双边计时}
\textbf{原则}：申论里说过的每个数字都要报得出作者和年份。先答是或否，再一句话守前提。家庭手势者的研究只有四个人，被问就大方承认。
\begin{qa}
\qarow{家庭手势者的研究有几个人？}{四个成年人，二〇一一年发在美国科学院院刊上。样本小，所以我方只用它说明一件事：有社会经验、没有语言，过了三就不准。}
\qarow{他们数得接近，是不是说明他们有数量感？}{是，大概多少人人都有，动物也有。缺的是精确，精确是跟着数词来的。}
\qarow{那门手语是谁造的？}{那群聋孩子。造出来之前，他们走不到；人造词，词带人走。}
\qarow{比划是不是也翻了一倍？}{是。守住，一倍；往前走，四倍。}
\qarow{失语的人还能不能思考？}{能，他们想的都是语言说得清的事，思考是有语言时搭起来的。}
\qarow{你方是不是把个人开除出了人类？}{不是。个人摸到的，是人类的起点；说出来、传下去，才是人类走到。}
\end{qa}
\end{stage}

\begin{stage}{7. 正三质询反二问题链}{90 秒，双边计时}
\textbf{总目标}：用反方自己的材料，拿到“一个人先摸到，人类要等词才走到”和“守住不等于往前走”。
\begin{chain}{链一：第二阶段}{爱因斯坦找词是为了告诉别人}
\q{爱因斯坦说，词要到第二阶段才去找，对吗？}{答“对”：下一问}
\q{找词，是为了能告诉别人，对吗？}{答“对”：下一问；答“不是”：请念原文“能传达给别人”}
\q{要是他一直没找到，人类知不知道相对论？}{答“不知道”：收束；答“知道”：请问怎么知道}
\cut{好，我换个问法：没找到词，论文写不写得出来？}
\close{谢谢。对方承认：一个人先摸到，人类要等他找到词，才走到。}
\end{chain}
\begin{chain}{链二：七十万年}{守住不等于往前走}
\q{您方说奥杜威石器没有像样的语言，也传了七十多万年，对吗？}{答“对”：下一问}
\q{这七十多万年里，石器有没有明显变好？}{答“没有”：下一问；答“有”：请报出处}
\q{那人类的边界，是守住的地方，还是走出去的地方？}{两种答法都收束到“往前走要靠传得准”}
\close{谢谢。七十万年守住了，也停住了。}
\end{chain}
\end{stage}

\begin{stage}{8. 正一对辩要点}{双方各 90 秒}
\textbf{目标}：收拢前两段，把“一个人摸到”和“人类走到”钉成全场的分歧。一辩的个人分主要在这里，发言短，追得紧。
\begin{points}{进攻点（4–5 个，每个 15–20 秒）}
\item 失语者、爱因斯坦、手艺人，都是一个人。请问，这个人走了以后，人类还剩下什么？
\item 对方说人先到、词后到。那词没来之前，这一步是在他一个人那里，还是在人类那里？
\item 对方说好用的工具不是边界。那请问，工具送不到的地方，下一代接不接得到？
\item 尼加拉瓜第一批人，二十五年的社会经验也替代不了语言。这是作者原话，对方承不承认？
\end{points}
\begin{points}{备用点（6–8 个）}
\item 对方说失语的人还能思考：我方一开始就承认，问题是他想的都是说得清的事。
\item 对方说默会知识：一个人知道的比说的多，人类知道的不比说的多。
\item 对方说汽车比走路快，公路不是边界：可人类的每一代，都只能从公路修到的地方出发。
\item 对方说言不尽意：那说不尽的一半，就留在了一个人心里。
\item 对方说维特根斯坦承认不可说：他让我们沉默，沉默的东西，没法接着往前推。
\item 对方说俄国人英国人都分得清蓝：看颜色是底子，给颜色分类、讲出光谱才是人类的边界。
\end{points}
\begin{points}{对方最可能问的三个问题 → 回应}
\item 你说不出的感受，是不是你的一部分？→ 是我的一部分，所以我说那是一个人的边界。
\item 爱因斯坦先想到相对论，那一刻人类有没有多一点东西？→ 多了一个人摸到的东西；写出来那天，人类才多了。
\item 说话翻四倍，比划翻一倍，那比划算不算走？→ 算守住。一倍和四倍，是原地和往前的差别。
\end{points}
\end{stage}

\begin{kit}{9. 正三申论}{套件 · 组装后 382–396 字 / 90 秒}
\assembly{开头 → 处理反方最强点（任选 1） → 推进论点二 → 收尾}
\begin{fixed}{开头}
谢谢主席。我先说这一篇做什么：先处理对方全场最有力的一点，再往前推我方第二点，说不出就接不上。
\end{fixed}
\begin{pick}{处理反方最强点}{1}
\begin{module}{如果对方全场最有力的是失语者还能思考}
对方全场最有力的，是失语的人还能思考。我方承认，能。可请评委想一想，这些人是怎么变成会算术、会下棋的人的？是在有语言的几十年里。做失语研究的学者自己也写了：就算语言在发育中是关键，成熟的大脑里也可能不再需要。这就是脚手架。真正的检验，是一辈子没有语言的人。我方举了两群，尼加拉瓜第一批手语者和家庭手势者，他们都走不到。
\end{module}
\begin{module}{如果对方全场最有力的是默会知识}
对方全场最有力的，是手艺、棋感这些说不出的本事。我方承认，老师傅手上有说不出的东西。可有一个问题绕不过去：师傅走了以后呢？说不出的那一点，徒弟只能自己从头摸。摸得到，是又一个人到了；摸不到，它就跟着师傅一起走了，再也找不回来。一代一代从零开始，人类就一直停在原地。守得住，却接不上，这就是我方说的边界。
\end{module}
\begin{module}{如果对方全场最有力的是石器没有语言也传了七十万年}
对方全场最有力的，是石器没有像样的语言，也传了七十多万年。我方承认，传了。可这七十多万年，出自同一个实验的论文：只看不交流的那组，一点没有进步；能说话的那组，翻了四倍。七十多万年，石器几乎没变样。人类不是一个只会守住的物种，是一个一代一代往前走的物种。守住，靠一双手；往前走，靠一张嘴。这组数据，站在我方这一边。
\end{module}
\begin{fallback}{对方最有力的一点不在以上几条时}
说白了，我方先把判准再说一遍。对方需要证明的，不是人有说不清的时候，而是在语言之外，人类能稳定地到达，还能留得住。一个人摸到了，那是他的边界；说出来、传下去，才是人类的。一个例子不管多动人，只要停在一个人、一次、一种说不清，就还在前一半。人类的边界，不看一个人摸到了哪里，要看下一个人接不接得住。
\end{fallback}
\end{pick}
\begin{fixed}{推进论点二}
我再往前推一步，讲一个词。一九七五年，美国康奈尔大学，一位叫卡米塔·伍德的职员，被上司骚扰了好几年，最后辞职。几个女人坐下来开会，想给这件事起个名字，定下了“性骚扰”。其中一位后来说：我们需要一个名字，让大家说的是同一件事。你想想，在那之前，无数人各自受着，各自说不清，各自以为只是自己倒霉。有了这个词，一个人的遭遇，才变成所有人都能认出来、能接着追究的事。后来的人不用再从头说起，接着往前走就行。
\end{fixed}
\begin{fixed}{收尾}
所以请评委记住今天的分歧：一个人摸到，是他的边界；说出来，传下去，下一个人接得住，人类才走到。
\end{fixed}
\end{kit}

\begin{stage}{10. 正三被质询防守预案}{反一质询，90 秒双边计时}
\textbf{原则}：性骚扰一词的每个细节都要报得出出处：History.com 2018 年的报道，学理是 Fricker 的诠释不正义。对方一定会说“她们早就说得清”，先认一半。
\begin{qa}
\qarow{一九七五年以前，她们知不知道自己受了委屈？}{知道自己难受，说得出一件件事；说不清这是同一种错，更没法让别人认出来。}
\qarow{这个名字是不是她们自己开会想出来的？}{是。人造词，词带人走；有了名字，无数人才认出是同一件事。}
\qarow{那是人先到，还是词先到？}{一个人先摸到，人类跟着词到。}
\qarow{失语的人还能不能思考？}{能。他们想的都是语言说得清的事，思考是有语言时搭起来的。}
\qarow{比划翻一倍，不说话也能教会吗？}{能守住，一倍；往前走，四倍。}
\qarow{爱因斯坦想到相对论那天，人类有没有多一点东西？}{多了一个人摸到的东西；他写出来那天，人类才多了。}
\end{qa}
\end{stage}

\begin{stage}{11. 正一质询反三问题链}{90 秒，双边计时}
\textbf{总目标}：全场最后一次质询，为结辩取材。拿到两个承认：说不出的那一点每一代从零摸；反方的感知例子证明的是底子，不是人类的边界。
\begin{chain}{链一：从零摸}{说不出的传不过去}
\q{您方说的棋感、手感，说不出来，对吗？}{答“对”：下一问}
\q{那师傅走了，徒弟没摸到，这一点还在不在？}{答“不在”：下一问；答“在”：请问在哪里}
\q{每一代都要从零去摸，算不算往前走？}{答“不算”：收束；答“算”：请问往前走了多少}
\close{谢谢。对方承认：说不出的那一点，每代从零开始。}
\end{chain}
\begin{chain}{链二：底子}{感知是人和动物共有的}
\q{英国人没有两个“蓝”字，也分得清，对吗？}{答“对”：下一问}
\q{分得清颜色，是不是人超出动物的那一截？}{答“不是”：下一问；答“是”：请问哪种动物分不清}
\q{那这个例子证明的，是人的底子，还是人类的边界？}{收束}
\close{谢谢。对方的例子，证明的是底子。}
\end{chain}
\end{stage}

\begin{stage}{12. 正二对辩要点}{双方各 90 秒}
\textbf{目标}：结辩前最后一次交锋，守住中立判准的两个字：稳定。把对方每一个例子都拉回“稳定地到达、留得住”这把尺子。
\begin{points}{进攻点（4–5 个，每个 15–20 秒）}
\item 判准是稳定地到达。一个说不清的瞬间，算不算稳定？
\item 说不出的手感，师傅一走就要从零摸。这算留得住吗？
\item 对方说快四倍只是工具好用。那没有这个工具，七十多万年，人类走了多远？
\item 失语的人还能思考，我方承认。请对方回答：他想的哪一件事，是语言说不清的？
\end{points}
\begin{points}{备用点（6–8 个）}
\item 对方说人先到词后到：一个人先到；人类跟着词到。
\item 对方说性骚扰早就说得清：说得清一件事，说不清是同一种错。
\item 对方说舍巴林还在作曲：乐谱是记号，曲子能传下去，靠的正是写下来的那部分。
\item 对方说婴儿会算：会的是三以内，动物也会。
\item 对方说造手语的是人：人造词，词带人走。
\item 对方说我方开除个人：个人摸到的，是人类的起点。
\end{points}
\begin{points}{对方最可能问的三个问题 → 回应}
\item 你方的“稳定”是不是专门给我方加的门槛？→ 不是，双方都要稳定：我方也不拿一个聪明的例子说语言什么都能。
\item 说不出的感受算不算人的一部分？→ 算一个人的；题目问的是人类。
\item 公路修到的地方就是人类的边界吗？→ 是下一代出发的地方；人类的边界，是一代接一代走出来的。
\end{points}
\end{stage}

\begin{kit}{13. 正一结辩}{套件 · 组装后 387–411 字 / 90 秒}
\assembly{归纳分歧（任选 1） → 处理最强点（任选 1） → 回应反一（任选 1） → 临场 → 论点收束 → 价值升华 → 结尾}
\begin{pick}{归纳分歧}{1}
\begin{module}{如果全场主战场落在失语者能不能思考}
谢谢主席。今天吵得最多的，是失语的人还能不能思考。我方一开始就承认，能。真正的分歧是：他想的那些，是不是语言到不了的地方。我方说，不是。
\end{module}
\begin{module}{如果全场主战场落在说不出的能不能传下去}
谢谢主席。今天吵得最多的，是说不出的东西传不传得下去。我方一开始就承认，能守住。真正的分歧是：守住，算不算人类在往前走。我方说，不算，接得上才算。
\end{module}
\begin{fallback}{主战场不清楚时}
谢谢主席。回到判准：语言到不了的地方，人类还能不能稳定地到达。今天的每一个例子，都要问一句：这是一个人摸到，还是人类走到？
\end{fallback}
\end{pick}
\begin{pick}{处理最强点}{1}
\begin{module}{如果反方全场最有力的是失语者还能思考}
反方最有力的，是失语者。我方承认，他们能想。可他们都是在有语言的几十年里，把思考一点点搭起来的。一辈子没有语言的人，尼加拉瓜那两群，一个答不对别人想错了什么，一个数不准三以上的东西。
\end{module}
\begin{module}{如果反方全场最有力的是默会知识}
反方最有力的，是手艺和棋感。我方承认，老师傅手上有说不出的东西。可师傅一走，徒弟只能从零去摸，一代一代从零开始。能守住，接不上；人类就一直停在原地。除了说出来，还有什么办法让下一代接着走？
\end{module}
\begin{fallback}{反方最有力的一点不在以上几条时}
我方承认，一个人的例子可以很动人。可只要停在一个人身上，一个人摸到，就只是他的边界；这个人走了以后，人类手里剩下的，只有他说出来、写下来的那部分，别的都跟着他走了。
\end{fallback}
\end{pick}
\begin{pick}{回应反一}{1}
\begin{module}{如果反一结辩收在语言只是工具不是边界}
反一说，好用的工具不是边界。可人类是一代接一代的人，每一代都只能从工具送到的地方出发。工具送不到的，下一代就接不到，只能从头再来。
\end{module}
\begin{module}{如果反一结辩收在说不出话的人还算不算完整的人}
反一问，说不出话的人算不算完整的人。算，我方从没说过不算。我方说的是，他心里没说出来的，别人接不到；所以我们更要早一点给人词。
\end{module}
\begin{fallback}{反一的收束不在以上几条时}
反一刚才的结辩，也请用同一把尺子量一量：它讲的，是一个人摸到了什么，还是人类稳定地走到了哪里。答案不一样，结论就不一样。
\end{fallback}
\end{pick}
\live{40}
\begin{fixed}{论点收束}
说到底，我方两点仍然成立。第一，没有词，一个人走不到人超出动物的那一截；第二，说不出，人类就接不上力。
\end{fixed}
\begin{fixed}{价值升华}
最后，请各位想一个孩子。五岁，聋，生在一个听得见的家里。这样的孩子很多。她爸妈很爱她，可是不会手语。她会笑，会生气，会指着饼干。她不笨。她只是还没有人，把词递到她手里。

我们今天争的，不是语言美不美。是一个人能走多远，取决于有没有人，早一点把词递给他。
\end{fixed}
\begin{fixed}{结尾}
请记住这句话：语言铺到哪里，人才走到哪里。谢谢。
\end{fixed}
\end{kit}
\begin{contingency}
\ifthen{时间不够}{先删“处理最强点”，再压“归纳分歧”；价值升华和结尾不删。}
\ifthen{反一结辩讲了一个具体的人}{临场位只说一句：他心里有的，我方不否认；他说出来的那部分，才留给了我们。}
\end{contingency}
\end{document}
